\chapter{CONCLUSION AND FUTURE RECOMMENDATIONS}

\section{Conclusion}
The \textbf{Personal Finance Manager (PFM)} project successfully demonstrates the potential of integrating Generative AI into everyday utility applications. By simplifying the expense logging process through Natural Language Processing, the system addresses the primary pain point of traditional finance apps---manual data entry friction. The use of a modern tech stack (React, FastAPI, Supabase) ensures the application is scalable, responsive, and maintainable.

The project met functionality objectives, including accurate parsing of expense text using Google Gemini, real-time data visualization, and secure user authentication. This system provides a robust foundation for personal financial management and showcases the practical application of Large Language Models (LLMs) in software engineering.

\section{Lesson Learnt/Outcome}
Through the development of this project, several key lessons were learned:
\begin{itemize}
    \item \textbf{AI Integration:} Learned how to effectively prompt engineer for consistent JSON outputs from LLMs.
    \item \textbf{State Management:} Gained deep understanding of React Context API for managing cross-component state.
    \item \textbf{Mobile Development:} Understood the nuances of building cross-platform apps using Capacitor.
    \item \textbf{API Design:} Learned importance of robust error handling in REST APIs.
\end{itemize}

\section{Future Recommendations}
While the current system is functional, several enhancements can be made in future iterations:

\begin{enumerate}
    \item \textbf{Bank SMS Integration:} Automate expense logging by reading transaction SMS messages (Android only).
    \item \textbf{OCR for Receipts:} Implement Optical Character Recognition (OCR) to allow users to scan physical receipts.
    \item \textbf{Budget Alerts:} Add push notifications when users exceed their set category budgets.
    \item \textbf{Investment Tracking:} Extend the system to track investments (Stocks, Mutual Funds) alongside daily expenses.
    \item \textbf{Offline Mode:} Enhance the PWA capabilities to allow offline logging, syncing when the internet is restored.
\end{enumerate}

% -------------------------------------------------------------------
% References
% -------------------------------------------------------------------
\renewcommand{\bibname}{References}
\begin{thebibliography}{9}

\bibitem{fastapi}
Tiangolo. (n.d.). \textit{FastAPI Framework}. [Online]. Available: \url{https://fastapi.tiangolo.com/}

\bibitem{react}
Meta Platforms. (n.d.). \textit{React - A JavaScript library for building user interfaces}. [Online]. Available: \url{https://reactjs.org/}

\bibitem{gemini}
Google. (2024). \textit{Gemini API Overview}. [Online]. Available: \url{https://ai.google.dev/}

\bibitem{supabase}
Supabase Inc. (n.d.). \textit{The Open Source Firebase Alternative}. [Online]. Available: \url{https://supabase.com/}

\bibitem{sommerville}
I. Sommerville, \textit{Software Engineering}, 9th ed. Boston, MA: Addison-Wesley, 2011.

\bibitem{pressman}
R. S. Pressman, \textit{Software Engineering: A Practitioner's Approach}, 8th ed. New York, NY: McGraw-Hill Education, 2014.

\end{thebibliography}
